\documentclass[../DoAn.tex]{subfiles}
\begin{document}

\section{Kết quả đạt được}

Đồ án xuất phát từ bài toán tối ưu phần tính toán của suy luận AI lượng tử hóa
trên bộ xử lý RISC-V nhúng. Trong bối cảnh các mô hình INT8 giúp giảm kích
thước dữ liệu nhưng vẫn tạo ra nhiều vòng lặp nhân--cộng, tải dữ liệu tuần tự
và điều khiển lặp, đồ án lựa chọn hướng tiếp cận ở mức tập lệnh. Mục tiêu của
hướng này là giảm số lệnh động và chi phí điều khiển trong các đoạn tính toán
lặp lại, nhưng vẫn giữ mô hình lập trình của một bộ xử lý RISC-V tổng quát thay
vì chuyển toàn bộ bài toán sang một bộ tăng tốc chuyên dụng.

Về mặt thiết kế, đồ án đã đặc tả một nhóm thao tác phục vụ suy luận INT8 gồm
\texttt{mac}, \texttt{dot4\_acc}, \texttt{relu}, \texttt{clamp}, \texttt{p.lw}
và \texttt{lp.setup}. Các thao tác này lần lượt đại diện cho nhân--cộng tích
lũy, tích vô hướng trên bốn phần tử INT8 đóng gói, kích hoạt, giới hạn miền giá
trị, tải dữ liệu kèm hậu tăng địa chỉ và thiết lập vòng lặp phần cứng. Trên cơ
sở đặc tả đó, đồ án đã xây dựng một luồng đánh giá gồm hai nhánh bổ trợ cho
nhau: nhánh mô phỏng hiện thực nhóm lệnh bằng mã hóa thử nghiệm trong vùng
\texttt{custom-0} của gem5, còn nhánh thực nghiệm ánh xạ cùng lớp ngữ nghĩa
sang các lệnh CORE-V/PULP trên lõi CV32E40P thông qua lớp firmware.

Đồ án cũng đã triển khai một luồng kiểm chứng hoàn chỉnh từ mô phỏng đến phần
cứng thật. Ở nhánh gem5, nhóm lệnh được đưa vào đường giải mã, lớp lệnh, ngữ
nghĩa thực thi, định dạng tải hậu tăng và trạng thái vòng lặp phần cứng; các
kiểm thử chức năng xác nhận những ngữ nghĩa chính hoạt động đúng trên các mô
hình CPU được chọn. Ở nhánh thực nghiệm, đồ án xây dựng SoC tối thiểu trên kit
DE2i-150/CV32E40P, xử lý ràng buộc thời gian để thiết kế đạt 50 MHz và đo hiệu
năng bằng các bộ đếm \texttt{mcycle}/\texttt{minstret}. Các kết quả tiêu biểu
cho thấy mức tăng tốc đạt 6,33 lần với biến thể dùng \texttt{dot4\_acc} và
13,10 lần khi kết hợp \texttt{dot4\_acc} với tải hậu tăng cùng vòng lặp phần
cứng; trong mô hình MNIST FC INT8 chạy bằng TensorFlow Lite Micro, đường tối ưu
khớp checksum và từng điểm đầu ra INT8 với đường tham chiếu, đồng thời đạt
12,39 lần ở khoảng đo có kiểm chứng và xấp xỉ 12,52 lần trong chế độ nhận ảnh
qua UART. Các kết quả này cho thấy hiệu quả của hướng mở rộng tập lệnh không
nằm ở việc gộp một phép toán riêng lẻ một cách cơ học, mà nằm ở việc rút gọn
đồng thời nhiều thành phần trong thân vòng lặp INT8, đặc biệt là tổ hợp tích vô
hướng đóng gói, tải hậu tăng và vòng lặp phần cứng.

Các kết quả trên cần được nhìn nhận cùng một số giới hạn của phạm vi đồ án. Thứ
nhất, mô hình gem5 đã được đồng bộ về ngữ nghĩa với nhánh phần cứng nhưng chưa
phải mô hình khớp từng chu kỳ với CV32E40P trên kit DE2i-150; do đó các kết quả
mô phỏng chỉ nên được hiểu như phân tích xu hướng. Thứ hai, đánh giá năng lượng
mới dựa trên VCD và Quartus Power Analyzer nên phù hợp để so sánh xu hướng tương
đối hơn là kết luận công suất tuyệt đối của kit. Cuối cùng, phạm vi mô hình thực
nghiệm hiện tập trung vào các nhân INT8 và mạng kết nối đầy đủ cho MNIST; các
mạng có cấu trúc tích chập hoặc các bài toán TinyML thực tế hơn chưa được khảo
sát trong phạm vi đồ án này.

\section{Hướng phát triển trong tương lai}

Từ các kết quả và hạn chế đã nêu, đồ án có thể được phát triển tiếp theo theo
ba hướng chính. Các hướng này tập trung vào việc làm cho mô hình đánh giá chính
xác hơn, lớp phần mềm thuận tiện hơn và phạm vi ứng dụng rộng hơn, thay vì mở
rộng dàn trải sang quá nhiều nhánh không trực tiếp nối với kết quả hiện tại.

Hướng phát triển thứ nhất là tiếp tục hoàn thiện nhánh mô phỏng. Mô hình gem5
có thể được điều chỉnh sâu hơn để phản ánh tốt hơn đặc điểm đường ống lệnh, độ
trễ lệnh và hành vi vòng lặp phần cứng của CV32E40P. Nếu có thêm đối chiếu với
mô phỏng RTL hoặc với các phép đo chi tiết hơn trên kit, nhánh gem5 sẽ không
chỉ đóng vai trò kiểm thử ngữ nghĩa mà còn hỗ trợ dự đoán thiết kế và phân tích
hiệu năng với độ tin cậy cao hơn.

Hướng phát triển thứ hai là cải thiện lớp công cụ phần mềm cho nhóm lệnh mở
rộng. Trong phiên bản hiện tại, firmware phát sinh các lệnh CORE-V/PULP thông
qua chỉ thị \texttt{.insn}, giúp kiểm soát chính xác mã lệnh nhưng chưa thuận
tiện cho lập trình và bảo trì lâu dài. Một bước tiếp theo phù hợp là bổ sung
hỗ trợ ở mức assembler, disassembler, intrinsic hoặc thư viện nhân tính toán chuẩn hóa
hơn, để các chương trình C/C++ và các nhân TensorFlow Lite Micro có thể gọi
nhóm thao tác tối ưu bằng giao diện rõ ràng thay vì phụ thuộc trực tiếp vào mã
hóa lệnh.

Hướng phát triển thứ ba là mở rộng phạm vi mô hình và toán tử được đánh giá.
Các kết quả hiện tại cho thấy nhóm thao tác đề xuất đặc biệt hiệu quả với lớp
kết nối đầy đủ có nhiều tích vô hướng INT8. Trong tương lai, cần tiếp tục thử
nghiệm với các mạng có lớp tích chập, tích chập theo chiều kênh
(depthwise convolution) hoặc các mô hình TinyML gần ứng dụng thực tế hơn.
Hướng này sẽ giúp đánh giá xem lợi ích của tích vô hướng đóng gói, tải hậu
tăng và vòng lặp phần cứng có duy trì khi dữ liệu có bố cục phức tạp hơn, kích
thước lớp không đều hơn và quá trình hậu xử lý lượng tử hóa chiếm tỷ trọng lớn
hơn hay không.

Tổng kết lại, đồ án đã hoàn thành một luồng nghiên cứu và thực nghiệm tương
đối đầy đủ cho hướng mở rộng tập lệnh RISC-V phục vụ suy luận AI lượng tử hóa:
từ phân tích bài toán, đặc tả nhóm thao tác, mô phỏng trong gem5, triển khai
SoC trên FPGA đến đánh giá mô hình suy luận thực tế bằng TensorFlow Lite Micro
trên phần cứng thật. Các kết quả thu được cho thấy hướng tiếp cận ở mức tập
lệnh có tiềm năng rõ ràng đối với các vòng lặp INT8 trên hệ nhúng, đồng thời
mở ra cơ sở để tiếp tục hoàn thiện công cụ, mô hình mô phỏng và phạm vi ứng
dụng trong các nghiên cứu sau.

\end{document}
